% =============================================================================
% Chapter 4: Method. Owner: Method writer. Plan: docs/thesis/REWRITE_PLAN.md, section 2 "4 Method".
% Numbers: numbers_method.tex (prefix \Method...). Tables: tables/T02, T07, T11. Figures: figures/tikz/F01, F03-F05.
% Trimmed 19 Oct (I/O table, HPH dimensions and ViT-S history numbers cut; see git 76aa1a4). Absorbs draft main.tex:216-580 (problem statement, architecture, loss
% equation, distillation mechanics, mixed training, schedule). Labels defined here: ch:method, sec:method-*,
% eq:method-*, fig:architecture, fig:anchor, fig:neck-chain, fig:distillation, tab:model-names,
% tab:loss-weights (T02), tab:metrics (T07), tab:student-recipe (T11).
% Labels used from other chapters: ch:data, ch:diagnostics, ch:results, app:reproduction.
% =============================================================================
\chapter{Method}
\label{ch:method}

This chapter describes the network, its losses, the training runs, the model names
(Section~\ref{sec:method-schedule}) and the evaluation. The two changes that matter most are in
Sections~\ref{sec:method-depth} and~\ref{sec:method-bonefix}: the training targets now describe the joint the network
anchors, and a constant rotation at the output turns four flipped bones back. Chapter~\ref{ch:diagnostics} shows how
the defects were found; Chapter~\ref{ch:results} measures what the changes did.

% -----------------------------------------------------------------------------
\section{Multi-HMR for Anny: task and architecture}
\label{sec:method-arch}

\subsection{The task}
\label{sec:method-task}

The input is one RGB image with camera intrinsics $\mathbf{K}$ (from an assumed field of view when they are unknown).
It shows an unknown number of people. For each person $i$ the network predicts the parent-relative rotations
$\mathbf{R}_i\in SO(3)^{\MethodNumBones}$ of Anny's bones, the phenotypes
$\boldsymbol{\phi}_i\in[0,1]^{\MethodNumPheno}$, and the image position $(u_i,v_i)$ and depth $z_i$ of one anchor
joint $a$, the head. The Anny body model $M$ \cite{bregier2025anny} turns $(\mathbf{R}_i,\boldsymbol{\phi}_i)$ into
vertices and joints, and the mesh is moved so that its anchor joint sits at the back-projected anchor pixel:
\begin{equation}
  \mathbf{V}_i = M_V(\mathbf{R}_i,\boldsymbol{\phi}_i) - M_{J,a}(\mathbf{R}_i,\boldsymbol{\phi}_i) + \mathbf{t}_i,
  \qquad
  \mathbf{t}_i = z_i\,\mathbf{K}^{-1}\,(u_i, v_i, 1)^{\top}.
  \label{eq:method-placement}
\end{equation}
The goal is a mesh with the right pose, shape and place in the camera frame, from one pass over the whole image.

\subsection{Architecture}
\label{sec:method-architecture}

The network is Multi-HMR \cite{baradel2024multihmr} with its output resized for Anny (Figure~\ref{fig:architecture}).
I did not change the backbone, the detection or the cross-attention head, so the architecture is no answer to
requirement 1b; the route to 1b is the smaller student of Section~\ref{sec:method-kd}.

\begin{figure}[!htb]
  \centering
  % F1 Architecture (Method 4.1). Owner: Method writer. tikzpicture only; the chapter wraps it in a figure.
% Dimensions: model.py:45-120 (backbone, detection MLPs, person_center head), :218-232 (HPH 1024 / 2 / 8 / 32),
% :421-427 (bone fix after rot6d_to_rotmat), :470-490 (Anny layer, placement at the anchor); numbers_method.tex.
% Libraries: positioning, arrows.meta, fit (thesis.tex loads them).
\begin{tikzpicture}[
  node distance=5mm and 4mm,
  box/.style={draw, rounded corners=2pt, align=center, font=\footnotesize, inner sep=3pt,
              minimum height=13mm},
  ours/.style={box, fill=teal!12, draw=teal!60!black},
  arr/.style={-{Stealth[length=2mm]}, semithick},
  tag/.style={font=\scriptsize\itshape, text=teal!50!black, align=left},
]
% row 1: image -> backbone -> detection
\node[box, text width=15mm] (img) {RGB image\\ $\MethodImgSize\times\MethodImgSize$\,px};
\node[box, text width=25mm, right=of img] (vit) {DINOv2 ViT-L/14\\ $\MethodTokenGrid\times\MethodTokenGrid$ tokens\\ width \MethodWidthL};
\node[box, text width=27mm, right=of vit] (det) {per token: detection score and 2D offset of the anchor joint (head)};
% row 2 (right to left): HPH -> Anny head -> bone correction
\node[box, text width=27mm, below=7mm of det] (hph) {HPH: cross-attention, \MethodHPHDepth{} blocks, \MethodHPHHeads{} heads, width \MethodHPHDim; one query per person};
\node[ours, text width=25mm, left=of hph] (out) {Anny head: $\MethodNumBones\times\MethodRotDims$ rotation values, \MethodNumPheno{} phenotypes, depth};
\node[ours, text width=15mm, left=of out] (fix) {bone correction\\ $R_b \leftarrow R_b C_b$\\ 0 parameters};
% row 3: body model
\node[ours, text width=44mm, below=6mm of out.south west, anchor=north] (body) {Anny body layer: mesh and joints, anchor joint placed at the predicted 2D position and depth};
% arrows
\draw[arr] (img) -- (vit);
\draw[arr] (vit) -- (det);
\draw[arr] (det) -- node[right, font=\scriptsize, align=left] {detected\\ tokens} (hph);
\draw[arr] (vit.south east) -- node[pos=0.45, above right, font=\scriptsize, inner sep=1pt] {all tokens} (hph.north west);
\draw[arr] (hph) -- (out);
\draw[arr] (out) -- (fix);
\draw[arr] (fix.south) -- (fix.south |- body.north);
% distillation box around the backbone
\node[draw, dashed, rounded corners=3pt, fit=(vit), inner sep=4pt,
      label={[font=\scriptsize, align=center]above:distillation: frozen ViT-L teacher,\\ ViT-B student with a $\MethodWidthB\to\MethodWidthL$ projector (Sec.~\ref{sec:method-kd})}] {};
% tags for what this thesis changed
\node[tag, right=2mm of det, text width=24mm] {ours: targets built at the same joint (Sec.~\ref{sec:method-depth})};
\node[tag, right=2mm of hph, text width=24mm] {ours: output resized for Anny, correction on four bones (Sec.~\ref{sec:method-bonefix})};
\node[tag, right=2mm of body, text width=40mm] {shaded = changed or added in this thesis; the rest is Multi-HMR};
\end{tikzpicture}

  \caption[The network]{The network. A $\MethodImgSize$-pixel image passes through a DINOv2 ViT-L backbone ($\MethodTokenGrid\times
  \MethodTokenGrid$ tokens). Every token gets a detection score and a 2D offset for the anchor joint, the head. For
  each detected person a cross-attention head (HPH) predicts $\MethodNumBones\times\MethodRotDims$ rotation values,
  \MethodNumPheno{} phenotypes and a depth; a constant correction turns four flipped bones back; the Anny layer builds
  the body and places it at the predicted distance. Shaded boxes are changed or added in this thesis; the dashed box
  marks the part that distillation replaces.}
  \label{fig:architecture}
\end{figure}

A DINOv2 ViT-L \cite{oquab2023dinov2} with weights from NAVER's Anny checkpoint turns the image into tokens; one MLP
scores each token for an anchor joint (the head, \texttt{-{}-person\_center head}) and a second predicts its 2D
offset. Each detected token becomes a query of the HPH, which lets the people attend to each other and to all image
tokens. The Anny head is what differs from the SMPL-X version: the rotations use the 6D representation
\cite{zhou2019continuity} with Gram--Schmidt, the phenotypes a sigmoid, and the depth Multi-HMR's focal-normalised
log-depth. Right after Gram--Schmidt four bones are multiplied by one constant rotation each
(Section~\ref{sec:method-bonefix}); then a fixed mask sets \MethodHelperBones{} helper bones (toes, breasts, face,
subdivision bones) to the identity, in the prediction and the ground truth alike, and keeps \MethodActiveBones{}. The
Anny layer runs inside the network and places the body with Equation~\eqref{eq:method-placement}, so the 3D and 2D
losses reach the network through it. The network has \MethodParamsTotalL\,M parameters (heads \MethodParamsHeadsL\,M)
and needs $\num{\MethodGflopsL}$ GFLOPs per image with one person.

% -----------------------------------------------------------------------------
\section{Losses in plain words}
\label{sec:method-losses}

Anny runs on each person's true parameters to give target vertices and joints; the anchor targets are built from one
joint (Section~\ref{sec:method-depth}). The total loss is
\begin{equation}
\begin{split}
\mathcal{L} = {}& \alpha_{\text{det}}\mathcal{L}_{\text{det}}
  + \alpha_{\text{off}}\mathcal{L}_{\text{off}}
  + \alpha_{\text{rot}}\mathcal{L}_{\text{rot}}
  + \alpha_{\text{shape}}\mathcal{L}_{\text{shape}} \\
  &+ \alpha_{\text{dist}}\mathcal{L}_{\text{dist}}
  + \alpha_{\text{t}}\mathcal{L}_{\text{t}}
  + \alpha_{\text{j3d}}\mathcal{L}_{\text{j3d}}
  + \alpha_{\text{v3d}}\mathcal{L}_{\text{v3d}} \\
  &+ \mathbb{1}[e \ge e_{\text{2D}}]\left(\alpha_{\text{j2d}}\mathcal{L}_{\text{j2d}}
  + \alpha_{\text{v2d}}\mathcal{L}_{\text{v2d}}\right)
  + \lambda_{\text{kd}}\mathcal{L}_{\text{kd}} + \lambda_{\text{kd-out}}\mathcal{L}_{\text{kd-out}},
\end{split}
\label{eq:method-loss}
\end{equation}
with $e$ the epoch; the distillation terms exist only for a student. Table~\ref{tab:loss-weights} says what each term
compares and gives its weights.

% T2 Loss terms and weights (Method 4.2). Owner: Method writer.
% Sources: defaults loss.py:436-452 (add_specific_args); terms loss.py:314-411 (forward);
% Final and Control: slurm/experiments/exp7_final_model.sh:210-222 (ARM fix / control), :759-771 (train.py command);
% Student: slurm/experiments/exp8_distill_vitb.sh header ("targets", "KD") and :319-331;
% distillation terms train.py:505-539; hand mask loss.py:231-253 (38 bones in the smoke log, H 0a step 6);
% shape mask loss.py:255-310 (off in all three columns; on in Final-masked, job 3612325).
% Boost renormalisation x4 -> x2.41 / x0.60: H 1 "x4 boost"; loss.py configure_joint_weights.
\begin{table}[!tbp]
\centering\footnotesize
\setlength{\tabcolsep}{3pt}
\begin{tabular}{@{}>{\raggedright}p{2.6cm}>{\raggedright}p{5.35cm}cccc@{}}
\toprule
Term (log name; \texttt{loss.py}) & What it compares, in plain words & Default & Final & Control & Student \tabularnewline
\midrule
Detection (\texttt{bce}; 314) & CornerNet focal loss on the per-token score ``an anchor joint lies in this token'' & \MethodWDet & \MethodWDet & \MethodWDet & \MethodWDet \tabularnewline
Offset (\texttt{offset}; 315) & L1 on the anchor's position inside its token & \MethodWOffset & \MethodWOffset & \MethodWOffset & \MethodWOffset \tabularnewline
Rotations (\texttt{rotmat}; 319--335) & L1 over the 9 entries of every bone's parent-relative rotation matrix & \MethodWRot & \MethodWRot & \MethodWRot & \MethodWRot \tabularnewline
\quad neck and hand boost & multiplier on neck, head and hand bones, then renormalised to mean 1 ($\times$\MethodBoostOn{} for them, $\times$\MethodBoostOff{} for the rest) & 1 & \MethodBoost & \MethodBoost & \MethodBoost \tabularnewline
Shape (\texttt{shape}; 336--343) & L1 on the \MethodNumPheno{} phenotypes after the sigmoid & \MethodWShapeDefault & \MethodWShape & \MethodWShape & \MethodWShape \tabularnewline
Distance (\texttt{dist}; 344) & L1 on the anchor's log-depth, normalised by the focal length & \MethodWDistDefault & \MethodWDist & \MethodWDistDefault & \MethodWDist \tabularnewline
Translation (\texttt{transl}; 345) & L1 on the anchor's 3D position in the camera frame & \MethodWDistDefault & \MethodWDist & \MethodWDistDefault & \MethodWDist \tabularnewline
3D joints (\texttt{j3d}; 367--377) & L1 on the \MethodNumBones{} joints, each skeleton moved to its own pelvis & \MethodWThreeD & \MethodWThreeD & \MethodWThreeD & \MethodWThreeD \tabularnewline
3D vertices (\texttt{v3d}; 378) & the same for the mesh vertices & \MethodWThreeD & \MethodWThreeD & \MethodWThreeD & \MethodWThreeD \tabularnewline
2D joints (\texttt{j2d}; 383--384) & L1 in pixels on the projected joints, placed by the predicted anchor and depth & \MethodWTwoDDefault & \MethodWTwoD & \MethodWTwoD & \MethodWTwoD \tabularnewline
2D vertices (\texttt{v2d}; 381--382) & the same for the projected vertices & \MethodWTwoDDefault & \MethodWTwoD & \MethodWTwoD & \MethodWTwoD \tabularnewline
\quad first epoch of the 2D terms & \texttt{-{}-start\_2d\_epoch} & \MethodStartTwoDDefault & \MethodStartTwoD & \MethodStartTwoD & \MethodStartTwoD \tabularnewline
Global rotation (\texttt{global\_rot}; 350--363) & L1 on the camera-frame orientation of head, neck03 and both wrists & 0 & 0 & 0 & 0 \tabularnewline
Feature distillation (\texttt{kd}) & KL divergence between the softened channel distributions of student and teacher tokens, $T=\MethodKDTemp$ & -- & -- & -- & \MethodKDLambda \tabularnewline
Output distillation (\texttt{kd\_out}) & L1 to the teacher's rotations, phenotypes and depth on the same people & -- & -- & -- & \MethodKDOutLambda \tabularnewline
\midrule
Joint the targets describe & \texttt{-{}-gt\_center\_joint}: transl, dist, detection token and offset targets & pelvis & head & pelvis & head \tabularnewline
Bone correction & constant output rotation on neck02, head, metacarpal1.L, finger4-1.R & none & on & none & on \tabularnewline
\bottomrule
\end{tabular}
\caption[Loss terms and their weights]{Loss terms and their weights. ``Default'' is the value in \texttt{loss.py}; Final, Control and Student are the
values of those runs (model names in Table~\ref{tab:model-names}). The total loss is the weighted sum of all rows above
the rule. For people from 3DPW the rotation and global-rotation terms skip the \MethodHandMaskBones{} finger, thumb
and metacarpal bones in all three runs; the mesh terms are not masked. The shape mask for 3DPW people is off in all
three runs and on only in Final-masked (Section~\ref{sec:method-real}). The Student keeps the teacher's bone correction
during training and is evaluated with its own, re-estimated after training (Section~\ref{sec:method-student}).}
\label{tab:loss-weights}
\end{table}


The detection term is the focal loss of CornerNet \cite{law2018cornernet} (the log calls it \texttt{bce}, but it is
not a binary cross-entropy). The neck and hand boost moves rotation loss between bones, it does not add any. The
shape L1 is smallest at the median of the plausible answers, so weak image cues pull predictions toward the training
median (Chapter~\ref{ch:diagnostics}). The 3D terms align both skeletons at the pelvis, so only distance and
translation see placement directly, and the 2D terms indirectly; before the depth fix the two direct terms made up
about \MethodLossShareDepth\% of the loss, against \MethodLossShareShape\% for shape. Training runs in fp32 with TF32
products (mixed precision gave NaN values with the Anny layer), and steps with NaN or Inf gradients are skipped
(\MethodFinalNaNSteps{} of $\num{\MethodFinalSteps}$ in Final). No weight was tuned in this thesis.

% -----------------------------------------------------------------------------
\section{Training schedule and model names}
\label{sec:method-schedule}

\paragraph{The Base lineage.} Training began from NAVER's public ViT-L checkpoint, Multi-HMR-Anny. Its head has a
different structure, so only the backbone was transferred. Four stages, L-stage 1 to 4, and a fifth run with shape
weight \MethodWShape{} led to \emph{Base} (epoch \MethodBaseEpoch): synthetic images only, and the start of every
fine-tune here. Continuations used a learning rate ten times below the cold start, which on a partly trained model
had caused hundreds of NaN gradient events per epoch.

\paragraph{Final and Control.} Final starts from Base and trains \MethodFinalEpochs{} epochs of
$\num{\MethodFinalIters}$ iterations at batch \MethodFinalBatch{} at a constant learning rate of \MethodFinalLR; the
last \MethodUnfrozenBlocks{} backbone blocks and the heads train. Three changes come on together: the depth fix with
distance and translation weight \MethodWDist{} (Section~\ref{sec:method-depth}), the bone correction
(Section~\ref{sec:method-bonefix}), and \MethodMixRatio\% real images per batch with the hand mask
(Section~\ref{sec:method-real}). The \MethodManifestImages{} images of the diagnostic manifest (Chapter~\ref{ch:data})
are never sampled. Training took $\num{\MethodFinalTrainSec}$\,s on one A100. Control is identical but keeps the
pelvis targets with weight \MethodWDistDefault{} and has no bone correction, so Control against Final isolates the two
fixes, and Base against Control shows the real images plus the restart cost. Each fix was tested alone only in short
pilots (Chapter~\ref{ch:diagnostics}).

\paragraph{Model names.} Table~\ref{tab:model-names} is the key to the names used from here on; each name is one run.
Table~\ref{tab:names} in Appendix~\ref{app:reproduction} gives each name's run, checkpoint, starting weights, data,
schedule and flags.

\begin{table}[!htb]
\centering\footnotesize
\setlength{\tabcolsep}{4pt}
\begin{tabular}{@{}>{\raggedright}p{2.75cm}>{\raggedright\arraybackslash}p{9.6cm}@{}}
\toprule
Name & Meaning \\
\midrule
Multi-HMR-Anny & NAVER's released ViT-L model for Anny; source of every backbone, baseline of requirement 1b \\
L-stage 1 to 4 & the transfer steps from Multi-HMR-Anny's backbone to Base \\
Base & end of the schedule, AnnyOne only; the start of every fine-tune \\
Mixed-L & Base fine-tuned with \MethodMixRatio\% 3DPW on the old labels; the previous best on 3DPW \\
Fresh-heads-L & fresh heads on Multi-HMR-Anny's backbone, \MethodMixRatio\% 3DPW (old labels) from the start \\
Control & Base, \MethodFinalEpochs{} epochs, \MethodMixRatio\% 3DPW (corrected labels), hand mask, neither fix \\
Final & Control plus the depth fix and the bone correction; the thesis model and the teacher \\
Final-masked & Final plus the shape mask on all \MethodShapeMaskCount{} converter constants \\
Final-partial-mask & Final plus the shape mask on the \MethodShapeMaskPartial{} constants other than age \\
Student-S/1 to /final & the ViT-S students of Base, before both fixes; history \\
Student & the ViT-B student distilled from Final; the answer to requirement 1b \\
\bottomrule
\end{tabular}
\caption[Model names used in this thesis]{Model names used in this thesis. Every name is one run; no run was repeated
with another seed. Full definitions: Table~\ref{tab:names} (Appendix~\ref{app:reproduction}).}
\label{tab:model-names}
\end{table}

% -----------------------------------------------------------------------------
\section{The depth anchor and its fix}
\label{sec:method-depth}

\paragraph{The mismatch.} The network places each person by the head: the detection token, the offset and the
predicted depth describe the head, and the mesh hangs from it (Equation~\eqref{eq:method-placement}). For AnnyOne
images, however, the training code built the translation, distance, token and offset targets at joint 0, the pelvis.
Every checkpoint before this thesis, Base included, was trained on these targets, so prediction and target described
two joints about half a body apart. A mesh hanging from a head placed at the pelvis is drawn too low; the 2D terms can
shrink this misfit by moving the body away, where it is drawn smaller, and neither the pelvis-aligned 3D terms nor the
weak direct depth terms object (Figure~\ref{fig:anchor}). The mismatch is certain; the mechanism is likely but was not
isolated. Section~\ref{sec:diag-depth} gives the evidence: all \MethodHoldoutPeople{} holdout people placed too far,
drawn at \MethodBaseSize{} of their size.

\begin{figure}[!htb]
  \centering
  % F3 The depth-anchor bug (Method 4.4 / Diagnostics 5.3). Owner: Method writer. tikzpicture only.
% Sources: model.py:399-403, :485-490 (mesh placed by the person_center joint = head); train.py:222-238 (targets at
% joint 0 unless --gt_center_joint person_center); H 3.1 (slope 0.33, intercept 2.36 m, size 0.63, 222 people all too
% far; Final 0.42 / 1.75 m); D slide 12; D notes 12 (3D height right). Numbers from numbers_method.tex.
\begin{tikzpicture}[
  font=\footnotesize,
  stick/.pic={
    \draw[pic actions] (0,1.8) circle (0.2);
    \draw[pic actions] (0,1.6) -- (0,0.8);
    \draw[pic actions] (-0.35,1.15) -- (0,1.45) -- (0.35,1.15);
    \draw[pic actions] (-0.3,0) -- (0,0.8) -- (0.3,0);
  },
  dot/.style={circle, fill, inner sep=1.5pt},
  arr/.style={-{Stealth[length=2mm]}, semithick},
  title/.style={font=\scriptsize\bfseries, anchor=north west, text width=3.8cm, align=left, inner sep=0pt},
  lab/.style={font=\scriptsize, align=left, text width=#1, inner sep=1pt},
]
% ---- panel A: what the model does
\node[title] at (0,3.1) {(a) Model: anchor at the head};
\pic[thick] at (0.5,0) {stick};
\node[dot, red!70!black] at (0.5,1.6) {};
\draw[red!70!black, thin] (0.65,1.6) -- (1.05,1.6);
\node[lab=2.6cm, anchor=west] at (1.05,1.35) {token, offset and depth describe the head joint; the mesh hangs from it};
% ---- panel B: what the targets said
\node[title] at (4.15,3.1) {(b) Targets before the fix: pelvis};
\pic[thick] at (4.65,0) {stick};
\node[dot, blue!70!black] at (4.65,0.8) {};
\draw[blue!70!black, thin] (4.8,0.8) -- (5.2,0.8);
\node[lab=2.6cm, anchor=west] at (5.2,1.1) {translation, distance, token and offset targets built at the pelvis};
% ---- panel C: side view, what training made of it
\node[title] at (8.3,3.1) {(c) Result: body pushed back};
\draw[thick, fill=gray!30] (8.3,0.8) rectangle (8.6,1.2);
\draw[thick, fill=gray!30] (8.6,0.9) -- (8.8,0.8) -- (8.8,1.2) -- (8.6,1.1) -- cycle;
\pic[dashed, gray] at (9.6,0) {stick};
\pic[thick] at (11.6,0) {stick};
\draw[gray, very thin] (8.8,1.0) -- (11.6,2.0);
\draw[gray, very thin] (8.8,1.0) -- (11.6,0.0);
\draw[arr] (9.9,2.35) -- node[above, font=\scriptsize] {2D terms} (11.3,2.35);
\node[font=\scriptsize, anchor=north] at (8.55,0.7) {camera};
\node[lab=1.2cm, align=center, anchor=north] at (9.6,-0.1) {true place};
\node[lab=1.9cm, align=center, anchor=north] at (11.6,-0.1) {same 3D height, drawn at \MethodBaseSize{} of its size};
% ---- bottom: predicted against true depth
\begin{scope}[shift={(1.0,-4.9)}, x=0.5cm, y=0.5cm]
  \draw[->] (0,0) -- (6.6,0);
  \draw[->] (0,0) -- (0,6.6);
  \foreach \t in {0,2,4,6} {
    \draw (\t,0.12) -- (\t,-0.12) node[below, font=\scriptsize] {\t};
    \draw (0.12,\t) -- (-0.12,\t) node[left, font=\scriptsize] {\t};
  }
  \node[font=\scriptsize, anchor=north] at (3.3,-0.9) {true pelvis depth (m)};
  \node[font=\scriptsize, rotate=90, anchor=south] at (-1.3,3.3) {predicted depth (m)};
  \draw[dashed, gray] (0,0) -- (6,6) node[right, font=\scriptsize] {ideal};
  \draw[very thick, red!70!black] (0.5,{\MethodBaseSlope*0.5+\MethodBaseIntercept}) -- (6,{\MethodBaseSlope*6+\MethodBaseIntercept});
  \draw[very thick, teal!60!black, densely dotted] (0.5,{\MethodFinalSlope*0.5+\MethodFinalIntercept}) -- (6,{\MethodFinalSlope*6+\MethodFinalIntercept});
\end{scope}
\node[lab=7.2cm, anchor=north west, text=red!70!black] at (5.4,-1.5)
  {Base (solid): predicted $= \MethodBaseSlope \times$ true $+\ \MethodBaseIntercept$\,m; all \MethodHoldoutPeople{} holdout people are placed too far.};
\node[lab=7.2cm, anchor=north west, text=teal!50!black] at (5.4,-2.4)
  {Final (dotted), after the fix: $\MethodFinalSlope \times$ true $+\ \MethodFinalIntercept$\,m; offset and drawn size moved, the slope stayed far from 1.};
\node[lab=7.2cm, anchor=north west] at (5.4,-3.3)
  {Why: the head joint is put where the pelvis target is, so the mesh hangs too low in the image; the 2D terms shrink that misfit by moving the body away.};
\end{tikzpicture}

  \caption[The depth-anchor bug]{The depth-anchor bug. (a) The model places each body by its head; (b) the training targets were built at
  the pelvis; (c) the 2D terms reconcile the two by moving the body away, where the same 3D body is drawn
  smaller. Bottom: the fitted line of predicted against true pelvis depth for the \MethodHoldoutPeople{}
  matched people of the AnnyOne holdout (setting V1), Base and Final. Schematic; the measured scatter is
  Figure~\ref{fig:depth-scatter}.}
  \label{fig:anchor}
\end{figure}

\paragraph{The fix.} A single training switch (\texttt{-{}-gt\_center\_joint person\_center}) builds all four anchor
targets at the joint the model anchors; the visibility filter and the one-person-per-token rule move with them. The
network does not change. I also raised the distance and translation weights from \MethodWDistDefault{} to
\MethodWDist, because they are the only direct depth terms; the pilots behind this choice are in
Table~\ref{tab:depth-runs}. Final uses the switch with weight \MethodWDist; Control keeps the pelvis targets with
weight \MethodWDistDefault.

% -----------------------------------------------------------------------------
\section{The bone correction}
\label{sec:method-bonefix}

\paragraph{What it corrects.} A per-bone dump on the holdout (Section~\ref{sec:diag-neck}) showed neck02 and the head
bone turned half a turn in all \MethodFlippedCkpts{} checkpoints of the Base lineage. The half-turns cancel, so the
head points the right way and PVE never showed the flip, but the neck skin between them is twisted to
\MethodBaseGirth{} of its true girth (Figure~\ref{fig:neck-chain}).

\begin{figure}[!htb]
  \centering
  % F4 The neck chain (Method 4.5 / Diagnostics 5.4). Owner: Method writer. tikzpicture only.
% Sources: H 3.2 (per-bone dump 3573518:12-15, :69: neck02 175.6, head 177.9, global head 12.0 deg, girth 0.46);
% 3579014:12, :37 (C_b 178.8 / 179.2 deg); 3579014:107 (girth 1.004 with the correction, no training);
% 3611336:198-207 (Final: girth 0.996); bone_fix.py docstring (R_b <- R_b C_b). Numbers from numbers_method.tex.
\begin{tikzpicture}[
  font=\footnotesize,
  bone/.style={draw, rounded corners=2pt, minimum width=15mm, minimum height=7mm, align=center},
  flip/.style={bone, fill=red!12, draw=red!60!black},
  okb/.style={bone, fill=teal!12, draw=teal!60!black},
  arr/.style={-{Stealth[length=2mm]}, semithick},
  turn/.style={-{Stealth[length=2mm]}, thick},
  note/.style={align=left, font=\scriptsize},
]
% ---- row 1: Base, before the correction
\node[font=\footnotesize\bfseries, anchor=west] at (-0.9,1.35) {Base, raw output};
\node[bone] (a1) at (0,0) {neck01};
\node[flip] (b1) at (2.4,0) {neck02};
\node[bone] (c1) at (4.8,0) {neck03};
\node[flip] (d1) at (7.2,0) {head};
\draw[arr] (a1) -- (b1);
\draw[arr] (b1) -- (c1);
\draw[arr] (c1) -- (d1);
\draw[turn, red!60!black] ([xshift=-4mm]b1.north) arc[start angle=180, end angle=0, radius=4mm];
\node[red!60!black, note, above=4mm of b1] {$\approx$\,180$^\circ$ twist (\MethodBaseNeckTwo$^\circ$)};
\draw[turn, red!60!black] ([xshift=4mm]d1.north) arc[start angle=0, end angle=180, radius=4mm];
\node[red!60!black, note, above=4mm of d1] {$\approx$\,180$^\circ$ back (\MethodBaseHead$^\circ$)};
\node[note, text width=3.0cm, anchor=west] at (8.3,0) {the two half-turns cancel: the head points right (\MethodBaseGlobalHead$^\circ$ error)};
\node[note, red!60!black, anchor=north] at (3.6,-0.55) {skin between neck02 and head twisted thin: girth \MethodBaseGirth{} of the true value};
% ---- row 2: with the correction
\node[font=\footnotesize\bfseries, anchor=west] at (-0.9,-1.75) {With the correction $R_b \leftarrow R_b C_b$};
\node[bone] (a2) at (0,-3.0) {neck01};
\node[okb] (b2) at (2.4,-3.0) {neck02\\[-1pt] {\scriptsize $\cdot\,C$ (\MethodCNeckTwo$^\circ$)}};
\node[bone] (c2) at (4.8,-3.0) {neck03};
\node[okb] (d2) at (7.2,-3.0) {head\\[-1pt] {\scriptsize $\cdot\,C$ (\MethodCHead$^\circ$)}};
\draw[arr] (a2) -- (b2);
\draw[arr] (b2) -- (c2);
\draw[arr] (c2) -- (d2);
\node[note, text width=3.0cm, anchor=west] at (8.3,-3.0) {both bones turned back by one constant each; 0 parameters};
\node[note, teal!50!black, anchor=north] at (3.6,-3.6) {girth \MethodFixGirth{} with no training; \MethodFinalGirth{} in Final};
\end{tikzpicture}

  \caption[The neck chain and its correction]{The neck chain neck01--neck02--neck03--head. Top: in Base, neck02 and the head bone are each turned about
  half a turn; the turns cancel, so the head points right, but the skin between them is twisted to \MethodBaseGirth{}
  of its true girth. Bottom: one constant rotation per bone turns both back. Girth on the AnnyOne holdout, setting V1.}
  \label{fig:neck-chain}
\end{figure}

\paragraph{How it works.} The parent-relative rotation of each flipped bone $b$ is multiplied on the right by a
constant, $\mathbf{R}_b\leftarrow\mathbf{R}_b\mathbf{C}_b$. If the prediction is $\hat{\mathbf{R}}_b\mathbf{F}_b$ for
the true rotation $\hat{\mathbf{R}}_b$ and a constant offset $\mathbf{F}_b$, then $\mathbf{C}_b=\mathbf{F}_b^{\top}$
restores the truth. I estimate $\mathbf{C}_b$ as the rotation closest to the mean of
$\mathbf{R}_{b,k}^{\top}\hat{\mathbf{R}}_{b,k}$ over the people $k$ of \MethodBoneFixImages{} random training images,
never the holdout. The script selects every bone whose median error exceeds \MethodBoneFixThr$^\circ$ with a constant
offset; for Base these are neck02, head, metacarpal1.L (base of the left index finger) and finger4-1.R (right ring
finger). The product comes before the helper mask, the kinematics and the losses, so training, evaluation and demos
all see corrected rotations. The constants are saved with the checkpoint; the correction adds no parameters. A
constant suffices because the HPH runs once from constant start values and never sees its own output.

\paragraph{Why not a loss, and limits.} The L1 distance over matrix entries has a barrier along a twist: at
135$^\circ$ it is larger than at 180$^\circ$ (\MethodLOneHalfTurn{} against \MethodLOneFullTurn), so a partial turn
back raises the loss; a geodesic loss and the neck boost did not undo the flip either (Section~\ref{sec:diag-neck}).
Without training the correction worsens the global head direction, because the head bone's offset varies between
people; Final therefore fine-tunes with the correction in place. A two-sided correction was not implemented, and the
origin of the flip is unknown.

% -----------------------------------------------------------------------------
\section{Training with real images}
\label{sec:method-real}

\paragraph{The mix.} A weighted sampler draws on average \MethodMixRatio\% of each batch from 3DPW-train converted to
Anny (Chapter~\ref{ch:data}; corrected labels, $\num{\MethodTrainThreeDPWImages}$ images) and \MethodSynthRatio\% from
AnnyOne; with more labelled people per AnnyOne image, roughly \MethodThreeDPWPeopleShare\% of the supervised people are
real (an estimate). The test split is never used for training. The mix answers requirement 2a.

\paragraph{The hand mask.} The converter fits only the root and the bones that carry a mapped SMPL joint, so every
finger, thumb and metacarpal bone of a 3DPW label is at rest. With \texttt{-{}-mask\_3dpw\_hands 1} the rotation terms
skip these \MethodHandMaskBones{} bones for 3DPW people, so the fingers are learned from AnnyOne only; the mesh terms
still carry the rest-pose hands.

\paragraph{The shape mask, tested.} The converter leaves \MethodShapeMaskCount{} phenotypes at constants (age
\MethodConstAge; muscle, weight, cup size, firmness and three ethnicity values at \MethodConstOther). With
\texttt{-{}-mask\_3dpw\_shape 1} the shape term ignores them for 3DPW people (Final-masked); the partial variant masks
only the \MethodShapeMaskPartial{} values other than age (Final-partial-mask). The mesh terms still carry the neutral
body. I kept the unmasked model as Final and teacher, because both masked models lost the gender gain on real photos
(Section~\ref{sec:res-mask}).

% -----------------------------------------------------------------------------
\section{Knowledge distillation}
\label{sec:method-kd}

A frozen teacher supervises a smaller student on the same image, on top of the task loss
(Figure~\ref{fig:distillation}).

\begin{figure}[!htb]
  \centering
  % F5 Distillation (Method 4.7). Owner: Method writer. tikzpicture only.
% Sources: slurm/experiments/exp8_distill_vitb.sh header (teacher frozen bf16, ViT-B student, --head_dim 1024 projector
% 768 -> 1024, heads copied and frozen 5 epochs, lambda_kd 1.0 T 4, lambda_kd_out 1.0); model.py:62-74 (feat_proj),
% :356-367 (KD tap on the tokens the heads see); train.py:505-539 (feature KL and output L1); distill.py (KL, T^2).
\begin{tikzpicture}[
  font=\footnotesize,
  node distance=4mm and 4mm,
  box/.style={draw, rounded corners=2pt, align=center, inner sep=3pt, minimum height=10mm},
  frozen/.style={box, fill=gray!15},
  ours/.style={box, fill=teal!12, draw=teal!60!black},
  arr/.style={-{Stealth[length=2mm]}, semithick},
  loss/.style={draw=red!60!black, dashed, {Stealth[length=2mm]}-{Stealth[length=2mm]}, semithick},
  lname/.style={font=\scriptsize, text=red!50!black, align=center, fill=white, inner sep=1pt},
]
\node[box, text width=10mm] (img) {image\\ $\MethodImgSize$\,px};
% teacher row
\node[frozen, text width=19mm, above right=4mm and 6mm of img] (tvit) {teacher ViT-L\\ frozen, bf16};
\node[frozen, text width=14mm, right=of tvit] (ttok) {tokens\\ width \MethodWidthL};
\node[frozen, text width=17mm, right=of ttok] (thead) {teacher heads\\ (frozen)};
\node[frozen, text width=14mm, right=of thead] (tout) {$R^{t}$, $\phi^{t}$, $d^{t}$};
% student row
\node[ours, text width=19mm, below right=4mm and 6mm of img] (svit) {student ViT-B\\ trained};
\node[ours, text width=14mm, right=of svit] (sproj) {projector\\ $\MethodWidthB\to\MethodWidthL$};
\node[ours, text width=17mm, right=of sproj] (shead) {heads copied from the teacher};
\node[ours, text width=14mm, right=of shead] (sout) {$R^{s}$, $\phi^{s}$, $d^{s}$};
% data flow
\draw[arr] (img) |- (tvit);
\draw[arr] (img) |- (svit);
\draw[arr] (tvit) -- (ttok);
\draw[arr] (ttok) -- (thead);
\draw[arr] (thead) -- (tout);
\draw[arr] (svit) -- (sproj);
\draw[arr] (sproj) -- (shead);
\draw[arr] (shead) -- (sout);
% losses
\draw[loss] (ttok) -- node[lname] {feature KL\\ $T=\MethodKDTemp$, $\lambda=\MethodKDLambda$} (sproj);
\draw[loss] (tout) -- node[lname] {L1 on rotations,\\ shape, depth; $\lambda=\MethodKDOutLambda$} (sout);
\draw[arr, densely dotted, gray] (thead.south) -- node[right, font=\scriptsize, text=gray, pos=0.5] {copy at start} (shead.north);
% labels and task loss
\node[font=\scriptsize, align=left, below=3mm of sout, text width=24mm] {plus the task loss of Table~\ref{tab:loss-weights} against the labels};
\node[font=\scriptsize, align=left, below=3mm of shead, text width=27mm] {frozen for the first \MethodStudentFreeze{} epochs};
\end{tikzpicture}

  \caption[Distillation of the Student]{Distillation of the Student. The frozen ViT-L teacher (Final, backbone in bf16) and the ViT-B student see
  the same image. A projector maps the student's tokens from width \MethodWidthB{} to \MethodWidthL, and heads copied
  from the teacher read them. The feature loss compares the projected tokens with the teacher's ($T=\MethodKDTemp$);
  the output loss compares rotations, phenotypes and depth of the same people. Shaded boxes train.}
  \label{fig:distillation}
\end{figure}

\paragraph{Feature distillation.} With teacher tokens $\mathbf{z}^{t}_n$ and the tokens $\mathbf{z}^{s}_n$ that the
student's heads see, both of width $C_t$,
\begin{equation}
\mathcal{L}_{\text{kd}} = T^2\cdot \frac{1}{N}\sum_{n=1}^{N}
  \mathrm{KL}\!\left(\sigma\!\left(\tfrac{\mathbf{z}^{t}_n}{T}\right)
  \,\Big\|\, \sigma\!\left(\tfrac{\mathbf{z}^{s}_n}{T}\right)\right),
\label{eq:method-kd}
\end{equation}
where $\sigma$ is the softmax over the channels of one token and $T$ the temperature \cite{hinton2015distilling}.
The loss sits between backbone and heads, so the student copies the teacher's representation, as in FitNets
\cite{romero2015fitnets}.

\paragraph{Head transfer and output distillation.} The student copies every teacher head tensor whose name and shape
match. Both run on the true person locations, so person $i$ of one is person $i$ of the other, and
\begin{equation}
\mathcal{L}_{\text{kd-out}} = \alpha_{\text{rot}}\left\|\mathbf{R}^{s}-\mathbf{R}^{t}\right\|_1
 + \alpha_{\text{shape}}\left\|\boldsymbol{\phi}^{s}-\boldsymbol{\phi}^{t}\right\|_1
 + \alpha_{\text{dist}}\left|d^{s}-d^{t}\right|
\label{eq:method-kdout}
\end{equation}
pulls the student toward the teacher's detached outputs, with the task weights of Table~\ref{tab:loss-weights}.

\paragraph{Warm-up freeze and projector-fed heads.} Copied heads that read through random input layers give confident
wrong outputs, and their gradients damage the copied weights; the copied tensors therefore stay frozen for the first
epochs. The final design removes the random input layers: a linear projector maps the student's tokens to the
teacher's width, every head is built at that width (\texttt{-{}-head\_dim}), all head tensors transfer, and
Equation~\eqref{eq:method-kd} trains the projector toward exactly the input the heads expect. TF32 products and the teacher's backbone in bfloat16 cut a ViT-S student's step from
\MethodItTimeBefore{} to \MethodItTimeAfter\,s on an A100, with unchanged evaluation numbers.

\subsection{History: the ViT-S students}
\label{sec:method-kd-history}

The first answer to the size requirement was a line of ViT-S students of Base, Student-S/1 to Student-S/final, trained
before both fixes and the label correction (numbers in Section~\ref{sec:res-history}). They taught three lessons.
Feature distillation alone leaves the student's heads unsupervised, so heads must be copied from the teacher. Copied
heads need projector-fed inputs, not random input layers. And the start matters most: a backbone already trained for
mesh recovery gave the largest gain, and the student was limited by the images it had seen, not by its size.

\subsection{The ViT-B student of Final}
\label{sec:method-student}

\paragraph{Recipe.} The Student follows the projector-fed design with a ViT-B backbone and the corrected training of
this chapter (Table~\ref{tab:student-recipe}). Its backbone starts from NAVER's Multi-HMR ViT-B, trained for mesh
recovery, and trains fully; the copied heads stay frozen for \MethodStudentFreeze{} epochs; the targets are exactly
the teacher's. I claim no augmentation: the job script names colour jitter, but no code reads that setting. The copied heads output the half-turn that the teacher's bone correction
undoes, so the student trains with that correction; afterwards its own is estimated on \MethodBoneFixImages{} training
images and used in every evaluation.

% T11 Student recipe and the 95% bars (Method 4.7.2; Results 6.10 refers to it). Owner: Method writer.
% Part A sources: slurm/experiments/exp8_distill_vitb.sh header (RECIPE: teacher, student, heads, bone fix, targets,
% KD, schedule) and :224 (LR_DECAY_EVERY = 2/5 of EPOCHS = 12), :319-331 (train.py flags);
% Final: exp7_final_model.sh:191-192, :224-225, :759-771. Cost: job 3613926 (exp8 smoke, A40) via H 0a; ViT-B
% backbone 86.58 M from model_cost.py (H 0a); projector 0.79 M and heads 14.0 M from the exp8 header.
% Part B sources: Final 3611336:191 (V1 PVE 73.7), :224-229 (3DPW PVE 89.0, MPJPE 88.3); Multi-HMR-Anny job 3614192
% (60.3 / 63.3 / 126.9 / 129.4) via H 0b; bars = error / 0.95 (exp8 header "THE 95% RULE"); 66.6 and 136.2 computed
% here from those values. Final holdout V1 MPJPE: not in any source read (the job re-measures it) -> TODO.
\begin{table}[!tbp]
\centering\footnotesize
\setlength{\tabcolsep}{4pt}
\begin{tabular}{@{}>{\raggedright}p{3.1cm}>{\raggedright}p{4.4cm}>{\raggedright\arraybackslash}p{4.6cm}@{}}
\toprule
\multicolumn{3}{@{}l}{\emph{A. Recipe}} \\
 & Final (the teacher) & Student \\
\midrule
Backbone & DINOv2 ViT-L/14, width \MethodWidthL & DINOv2 ViT-B/14, width \MethodWidthB \\
Backbone initialisation & Base (from NAVER's Anny checkpoint) & NAVER's Multi-HMR ViT-B, trained for mesh recovery \\
Trainable during training & last \MethodUnfrozenBlocks{} of \MethodBlocksL{} blocks and the heads; frozen when teaching (bf16) & whole backbone, projector; heads after epoch \MethodStudentFreeze \\
Heads & HPH at width \MethodHPHDim, \MethodParamsHeadsL\,M & copied from the teacher, fed by a projector \MethodWidthB{}$\to$\MethodWidthL{} (\MethodStudentProjParams\,M) \\
Parameters & \MethodParamsTotalL\,M & \MethodStudentParams\,M (\MethodRatioParams$\times$ fewer) \\
GFLOPs (672\,px, one person) & \num{\MethodGflopsL} & \MethodStudentGflops{} (\MethodRatioGflops$\times$ fewer) \\
ms per image (A40, batch 1) & \MethodMsL & \MethodStudentMs{} (\MethodRatioMs$\times$ faster) \\
Schedule & \MethodFinalEpochs{} epochs $\times$ \num{\MethodFinalIters} iterations $\times$ batch \MethodFinalBatch{} $=$ \num{\MethodFinalImages} images & \MethodStudentEpochs{} epochs $\times$ \num{\MethodStudentIters} iterations $\times$ batch \MethodStudentBatch{} $=$ \num{\MethodStudentImages} images \\
Learning rate & \MethodFinalLR, constant & \MethodStudentLR, halved every \MethodStudentLRDecay{} epochs \\
Task targets & head anchor, distance and translation $\times$\MethodWDist, \MethodMixRatio\% 3DPW-train (v23), hand mask, manifest excluded & identical to the teacher \\
Distillation & -- & feature KL $\lambda=\MethodKDLambda$, $T=\MethodKDTemp$; output L1 $\lambda=\MethodKDOutLambda$ \\
Bone correction & four bones, estimated on Base & teacher's during training; own one re-estimated on \MethodBoneFixImages{} training images afterwards \\
\midrule
\multicolumn{3}{@{}l}{\emph{B. The 95\% bars: student error $\le$ reference error / \MethodRuleFactor}} \\
 & against Final & against Multi-HMR-Anny \\
\midrule
AnnyOne holdout PVE (V1) & \MethodTeacherHoldoutPVE{} $\to$ \textbf{\MethodBarHoldout} & \MethodNaverHoldoutPVE{} $\to$ \MethodBarNaverHoldout{} (likely in-sample) \\
AnnyOne holdout MPJPE (V1) & \TODO{Final V1 MPJPE, re-measured in the student job} & \MethodNaverHoldoutMPJPE{} $\to$ \MethodBarNaverHoldoutMPJPE \\
3DPW test PVE (v23) & \MethodTeacherThreeDPWPVE{} $\to$ \textbf{\MethodBarThreeDPW} & \MethodNaverThreeDPWPVE{} $\to$ \MethodBarNaverThreeDPW \\
3DPW test MPJPE (v23) & \MethodTeacherThreeDPWMPJPE{} $\to$ \MethodBarThreeDPWMPJPE & \MethodNaverThreeDPWMPJPE{} $\to$ \MethodBarNaverThreeDPWMPJPE \\
\bottomrule
\end{tabular}
\caption{The Student (ViT-B) next to its teacher Final. Part A: recipe and measured cost. Part B: the error bars of
the 95\% rule in mm; the job computes the bars from the teacher re-measured in the same job, so these values are for
orientation. The rule also asks for fewer parameters, fewer FLOPs and a faster forward pass. Read as ``error at most
1.05 times the teacher's'', the bars against Final would be \MethodBarHoldoutAlt{} and \MethodBarThreeDPWAlt\,mm. The
Student's own errors are in Chapter~\ref{ch:results}.}
\label{tab:student-recipe}
\end{table}


\paragraph{The 95\% rule.} Requirement 1b asks for a smaller model within 95\% of the large one. I read it as: student
error at most the teacher's error divided by \MethodRuleFactor, for PVE and MPJPE on the AnnyOne holdout (V1) and the
3DPW test set, with fewer parameters, fewer FLOPs and a faster forward pass (bars in Table~\ref{tab:student-recipe}).
NAVER's model was trained on the AnnyOne release, which includes our \MethodHoldoutN{} holdout images, so its
\MethodNaverHoldoutPVE\,mm is most likely in-sample; a student of a \MethodTeacherHoldoutPVE\,mm teacher cannot be
expected to reach its bar of \MethodBarNaverHoldout\,mm, while its 3DPW bar of \MethodBarNaverThreeDPW\,mm lies above
the teacher's own error. Which reference counts is the supervisor's
decision; Chapter~\ref{ch:results} reports the student against both.

% -----------------------------------------------------------------------------
\section{Evaluation protocol}
\label{sec:method-eval}

Table~\ref{tab:metrics} lists the metrics. Errors are means over matched people: a greedy 2D matcher pairs each
labelled person with at most one detection, and recall is the share of labelled people that got one.

% T7 Metrics (Method 4.8). Owner: Method writer.
% Sources: docs/handover.md section 1 "Metrics" and "Terms" (trace metrics, girth, PVE_neck, verdicts);
% docs/deck_explained_2026-10-19.md part 2 and slide 9; scripts/diagnostics/diag_holdout_trace.py:14-24 (settings);
% scripts/diagnostics/eval_relative_human.py docstring (AUC, gender accuracy, IoU, child threshold);
% README :321 (V0 - V1 = +3.8 mm); REWRITE_PLAN.md 1.1 #15 (163-joint MPJPE, not comparable).
\begin{table}[!htb]
\centering\footnotesize
\setlength{\tabcolsep}{3pt}
\begin{tabular}{@{}>{\raggedright}p{2.6cm}>{\raggedright}p{6.4cm}>{\raggedright\arraybackslash}p{3.4cm}@{}}
\toprule
Metric & What it measures & Blind to \\
\midrule
PVE (mm) & vertex error after both pelvises are put at the same point & placement, depth, size \\
PA-PVE (mm) & PVE after the best rotation, scale and shift (Procrustes) & also global orientation \\
MPJPE (mm) & PVE on the \MethodNumBones{} Anny joints & as PVE \\
Recall, precision (\%) & labelled people matched; detections matched & false detections; bystanders \\
Root error (m) & 3D pelvis distance, no alignment & pose and shape \\
Size ratio & drawn head-to-pelvis distance, predicted over true (1 is right) & 3D body size \\
Depth slope / intercept & line fitted to predicted against true pelvis depth (1 / 0\,m) & errors around the line \\
Bone angle error (deg) & angle of $\hat{R}_b^{\top}R_b$ for a parent-relative rotation; ${\approx}180$ = flipped & compensating bones \\
Head direction (deg) & camera-frame head orientation error & the chain below the head \\
Neck girth ratio & neck-skin radius, predicted over true (1 is right) & neck length \\
PVE\textsubscript{neck} (mm) & vertex error over the neck skin, centred at neck01 & the rest of the body \\
Child AUC & chance that a child gets a lower predicted age than an adult (0.5 = chance) & undetected people \\
Gender accuracy (\%) & predicted gender on the right side of 0.5 (50 = chance) & undetected people \\
\bottomrule
\end{tabular}
\caption[Metrics]{Metrics used in this thesis. All errors are means over matched people; the last two are on
Relative Human; MPJPE is not comparable with published 3DPW numbers. On the AnnyOne holdout the headline
setting is V1 (helper bones masked); V0 leaves them unmasked and reads \MethodVzeroGap\,mm higher in PVE on the same
people. 3DPW numbers are against our converted labels only (Chapter~\ref{ch:data}).}
\label{tab:metrics}
\end{table}


\paragraph{AnnyOne holdout.} The last \MethodHoldoutN{} AnnyOne images (indices $\num{\MethodHoldoutFirst}$ to
$\num{\MethodHoldoutLast}$), never trained on and without children (Chapter~\ref{ch:data}). The headline setting V1
masks the helper bones, as the demos do; V0 does not. A trace script adds the placement and per-bone metrics.

\paragraph{3DPW test.} All $\num{\MethodTestImages}$ test images ($\num{\MethodTestPeople}$ people) of the corrected
labels, same script for every model, setting V0. MPJPE averages all \MethodNumBones{} Anny joints and the conversion
itself errs by about \MethodLabelFloor\,mm (Chapter~\ref{ch:data}), so these numbers are not comparable with published
3DPW numbers.

\paragraph{Relative Human and demo images.} Relative Human gives real photos with an age group, a gender and a box per
person. Its validation split (\MethodRHImages{} images, $\num{\MethodRHPeople}$ people) serves for
evaluation only: boxes from the projected vertices, matching at IoU $\ge$ \MethodRHIoU, child-shaped meaning a
predicted age below \MethodChildThr. The \MethodDemoPhotos{} demo photos and \MethodDemoFrames{} 3DPW test frames are
rendered with each body placed from its detection box, which hides depth errors, and are judged by eye.

\paragraph{One run, no seed.} The training code sets no random seed, and every model was trained once. Two identical
pilots show that one run resolves about \MethodResolvePVE\,mm PVE, \MethodResolveHead$^\circ$ head direction and
\MethodNoiseRecall{} recall points (Section~\ref{sec:diag-noise}); smaller gaps are within run-to-run variation.
Restarting a fine-tune of Base with nothing changed already costs about \MethodRestartCost\,mm PVE, so every fine-tune
is compared with its own control. The pelvis error has no measured noise floor, and no bootstrap intervals were
computed.
